\begin{abstract}
A local routing defense can identify compromised nodes while leaving legitimate destinations unreachable. This study examines that distinction under concurrent DODAG Information Solicitation (DIS) flooding and incorrect-parent Destination Advertisement Object (DAO) manipulation in RPL networks. Adaptive Evidence Fusion for RPL (AEF-RPL) extends an independent DIS-FDM implementation with calibrated first-hop DAO-consistency evidence, persistent evidence fusion, selective DAO handling, and a tabular Q-learning decision layer evaluated against a learning-disabled variant. A controlled Contiki-NG/Cooja evaluation comprises 1960 archived 30-node runs, twenty fresh-seed benign mobility runs, and a 740-run extension covering attack rates, component ablations, reception loss, temporal alerts, host callback cost, and scale feasibility. In the default-rate extension, mean downward delivery is 95.73\% for AEF learning versus 78.70\% for DIS-FDM in static networks, but broad multiplicity-adjusted tests do not establish superiority. Mobility exposes substantial false isolation, and the benign learning comparison does not demonstrate a false-positive reduction. Ablations associate selective DAO repair with service preservation more strongly than learning with reliable attribution. Attack-free 500/1000-node controls fail connectivity requirements, preventing a validated large-scale defense claim. Independently reparsed logs reproduce the completed extension's numerical results. The findings establish an implementation-specific service--attribution trade-off and identify the validation needed before deployment claims.
\end{abstract}
\begin{IEEEkeywords}
RPL security, intrusion detection, DAO manipulation, DIS flooding, Dempster--Shafer evidence fusion, Q-learning, reproducible simulation.
\end{IEEEkeywords}
